% tex/impl/commands/numbers.tex
\ProvidesFile{numbers.tex}[2026/08/24 Number-system notation]

% Integer construction from equivalence classes of pairs of natural numbers.
\DeclareRobustCommand{\IntPairRel}{%
  \ensuremath{\mathrel{\sim_{\mathbb{Z}}}}%
}
\DeclareRobustCommand{\IntClass}[2]{%
  \ensuremath{\left[#1,#2\right]_{\mathbb{Z}}}%
}
\DeclareRobustCommand{\IntEmbed}{%
  \ensuremath{\iota_{\mathbb{Z}}}%
}
\DeclareRobustCommand{\NatCopyInInt}{%
  \ensuremath{\mathbb N_{\mathbb Z}}%
}
% Typed integer numerals.  The visible glyphs deliberately agree with the
% natural-number glyphs; \IntOfNat keeps the embedding explicit whenever the
% ambient carrier does not determine the type unambiguously.
\DeclareRobustCommand{\IntNumeral}[1]{%
  \ensuremath{\mathord{#1}}%
}
\DeclareRobustCommand{\IntZero}{\IntNumeral{0}}
\DeclareRobustCommand{\IntOne}{\IntNumeral{1}}
\DeclareRobustCommand{\IntOfNat}[1]{%
  \ensuremath{\IntEmbed\!\left(#1\right)}%
}
\DeclareRobustCommand{\IntSucc}{%
  \ensuremath{S_{\mathbb{Z}}}%
}
\DeclareRobustCommand{\IntPred}{%
  \ensuremath{P_{\mathbb{Z}}}%
}
\DeclareRobustCommand{\IntegerPeano}[1]{%
  \ensuremath{\mathsf{ZPeano}\!\left(#1\right)}%
}
\DeclareRobustCommand{\IntegerArithmetic}[1]{%
  \ensuremath{\mathsf{ZArith}\!\left(#1\right)}%
}

% Rational construction from equivalence classes of integer fractions.
\DeclareRobustCommand{\PositiveIntegers}{%
  \ensuremath{\mathbb Z_{>0}}%
}
\DeclareRobustCommand{\RatPairCarrier}{%
  \ensuremath{D_{\mathbb Q}}%
}
\DeclareRobustCommand{\RatPairRel}{%
  \ensuremath{\mathrel{\sim_{\mathbb Q}}}%
}
\DeclareRobustCommand{\RatClass}[2]{%
  \ensuremath{\left[#1,#2\right]_{\mathbb Q}}%
}
\DeclareRobustCommand{\RatEmbed}{%
  \ensuremath{\iota_{\mathbb Q}}%
}
\DeclareRobustCommand{\NatCopyInRat}{%
  \ensuremath{\mathbb N_{\mathbb Q}}%
}
\DeclareRobustCommand{\IntCopyInRat}{%
  \ensuremath{\mathbb Z_{\mathbb Q}}%
}
% Typed rational numerals.  As for the integers, the visible glyph is shared
% with the natural-number notation; the explicit embedding macros remain
% available whenever the ambient carrier does not determine the type.
\DeclareRobustCommand{\RatNumeral}[1]{%
  \ensuremath{\mathord{#1}}%
}
\DeclareRobustCommand{\RatZero}{\RatNumeral{0}}
\DeclareRobustCommand{\RatOne}{\RatNumeral{1}}
\DeclareRobustCommand{\RatOfInt}[1]{%
  \ensuremath{\RatEmbed\!\left(#1\right)}%
}
\DeclareRobustCommand{\RatOfNat}[1]{%
  \ensuremath{\RatOfInt{\IntOfNat{#1}}}%
}
\DeclareRobustCommand{\RatOrderGraph}{%
  \ensuremath{R^{\leq}_{\mathbb Q}}%
}
\DeclareRobustCommand{\RatRecipRep}[2]{%
  \ensuremath{\operatorname{rec}_{\mathbb Q}\!\left(#1,#2\right)}%
}
\DeclareRobustCommand{\RationalOrderedField}[1]{%
  \ensuremath{\mathsf{QOF}\!\left(#1\right)}%
}

% Dedekind-cut model of the real numbers.
\DeclareRobustCommand{\RealCut}[1]{%
  \ensuremath{\operatorname{Cut}_{\mathbb Q}\!\left(#1\right)}%
}
\DeclareRobustCommand{\RealHat}[1]{\ensuremath{\widehat{#1}}}
\DeclareRobustCommand{\RealEmb}{\ensuremath{\iota_{\mathbb R}}}
\DeclareRobustCommand{\NatCopyInReal}{%
  \ensuremath{\mathbb N_{\mathbb R}}%
}
\DeclareRobustCommand{\IntCopyInReal}{%
  \ensuremath{\mathbb Z_{\mathbb R}}%
}
\DeclareRobustCommand{\RatCopyInReal}{%
  \ensuremath{\mathbb Q_{\mathbb R}}%
}
\DeclareRobustCommand{\RealLe}{\ensuremath{\mathrel{\leq_{\mathbb R}}}}
\DeclareRobustCommand{\RealLt}{\ensuremath{\mathrel{<_{\mathbb R}}}}
\DeclareRobustCommand{\RealZero}{\ensuremath{0_{\mathbb R}}}
\DeclareRobustCommand{\RealOne}{\ensuremath{1_{\mathbb R}}}
\DeclareRobustCommand{\RealAdd}{\ensuremath{\mathbin{\oplus}}}
\DeclareRobustCommand{\RealMul}{\ensuremath{\mathbin{\odot}}}
\DeclareRobustCommand{\RealPosMul}{\ensuremath{\mathbin{\odot_{+}}}}
\DeclareRobustCommand{\RealNeg}[1]{\ensuremath{\mathop{\ominus}\!#1}}
\DeclareRobustCommand{\RealPosInv}[1]{%
  \ensuremath{\operatorname{inv}_{+}\!\left(#1\right)}%
}
\DeclareRobustCommand{\RealInv}[1]{%
  \ensuremath{\operatorname{inv}_{\mathbb R}\!\left(#1\right)}%
}
\DeclareRobustCommand{\RealAbs}[1]{%
  \ensuremath{\left\lvert #1\right\rvert_{\mathbb R}}%
}
\DeclareRobustCommand{\RealDist}[2]{%
  \ensuremath{d_{\mathbb R}\!\left(#1,#2\right)}%
}
\DeclareRobustCommand{\CompleteOrderedField}[1]{%
  \ensuremath{\mathsf{COF}\!\left(#1\right)}%
}

% Bases and digit sets.
\DeclareRobustCommand{\BaseOK}[1]{%
  \ensuremath{\mathsf{Bas}\!\left(#1\right)}%
}
\DeclareRobustCommand{\DigitSet}[1]{%
  \ensuremath{\mathcal{D}_{#1}}%
}
\DeclareRobustCommand{\DecDigitSet}{%
  \ensuremath{\mathcal{D}}%
}
\DeclareRobustCommand{\BinDigitSet}{%
  \ensuremath{\DigitSet{2}}%
}

% Digit-sequence predicates. The last argument is the highest index.
\DeclareRobustCommand{\DigitSeq}[3]{%
  \ensuremath{\mathsf{Ziff}_{#1}\!\left(#2,#3\right)}%
}
\DeclareRobustCommand{\DecDigitSeq}[2]{%
  \ensuremath{\mathsf{Ziff}\!\left(#1,#2\right)}%
}
\DeclareRobustCommand{\BinDigitSeq}[2]{%
  \ensuremath{\mathsf{Ziff}_{2}\!\left(#1,#2\right)}%
}
\DeclareRobustCommand{\NormDigitSeq}[3]{%
  \ensuremath{\mathsf{NZiff}_{#1}\!\left(#2,#3\right)}%
}

% Peano initial segment {i in N | i < n}.
\DeclareRobustCommand{\NatSeg}[1]{%
  \ensuremath{\mathbb{N}_{<#1}}%
}
\DeclareRobustCommand{\PosNatSeg}[1]{%
  \ensuremath{\mathbb{N}_{1,\ldots,#1}}%
}

% Place-value term sequence and evaluation. The last argument is the length of
% the segment.
\DeclareRobustCommand{\PlaceTermSeq}[2]{%
  \ensuremath{w_{#1,#2}}%
}
\DeclareRobustCommand{\tPlaceTermSeq}[2]{%
  \ensuremath{t_{#1,#2}}%
}
\DeclareRobustCommand{\PlaceVal}[3]{%
  \ensuremath{\operatorname{val}_{#1}\!\left(#2,#3\right)}%
}
\DeclareRobustCommand{\DecPlaceVal}[2]{%
  \ensuremath{\operatorname{val}\!\left(#1,#2\right)}%
}

% Visible digit blocks and numerals.
\newcommand{\DigitBlockRaw}[2]{#1_{#2}#1_{#2-1}\cdots #1_0}
\DeclareRobustCommand{\BaseNum}[3]{%
  \ensuremath{\mathord{\DigitBlockRaw{#2}{#3}}_{#1}}%
}
\DeclareRobustCommand{\BinNum}[2]{%
  \ensuremath{\mathord{\DigitBlockRaw{#1}{#2}}_{2}}%
}
\DeclareRobustCommand{\DecNum}[2]{%
  \ensuremath{\DigitBlockRaw{#1}{#2}}%
}

% Literal binary digit strings, used for examples.
\DeclareRobustCommand{\BinDigitString}[1]{%
  \ensuremath{\mathord{#1}_{2}}%
}
